% ============================================================================
% Section 4: L_RAG System Design and Methodology
% ============================================================================
\section{L\_RAG System Design and Methodology}
\label{sec:system-design}

L\_RAG turns the structured corpus and benchmark into an evidence-oriented
retrieval and generation pipeline. Offline processing creates normalized
records, graph artifacts, and vector or sparse indexes. Online processing
resolves query scope, retrieves candidate chunks, optionally expands or
filters them with the graph, reranks a bounded pool, and serializes the final
context for a citation-grounded generator. The same document, provision, and
chunk identifiers introduced in Section~\ref{sec:dataset-construction} are
carried through every stage.

\subsection{End-to-End System Architecture}
\label{sec:architecture}

The architecture has five operational phases: (1) query and temporal-scope
resolution, (2) dense and optional sparse retrieval, (3) graph filtering or
context expansion, (4) rank fusion and candidate reranking, and (5) prompt
construction, generation, and cited output. \hyperref[fig:architecture]{Figure~\ref*{fig:architecture}}
shows the data flow and the point at which each family of ablations changes a
component.

\begin{figure*}[t]
\centering
\resizebox{\textwidth}{!}{%
\begin{tikzpicture}[node distance=1.45cm, auto,
    block/.style={rectangle, draw=blue!80, fill=blue!5, thick, minimum width=3cm, minimum height=1cm, align=center, rounded corners},
    kgblock/.style={rectangle, draw=green!60!black, fill=green!5, thick, minimum width=3cm, minimum height=1cm, align=center, rounded corners},
    llmblock/.style={rectangle, draw=red!70, fill=red!5, thick, minimum width=3cm, minimum height=1cm, align=center, rounded corners},
    line/.style={draw, -latex', thick}]
    \node [block] (query) {User query $q$\\+(date/scope)};
    \node [block, right=1.3cm of query] (retrieve) {Dense E5\\+ optional BM25};
    \node [kgblock, below=1.15cm of retrieve] (graph) {G-LRAG overlay\\filter/expand};
    \node [block, right=1.3cm of retrieve] (fuse) {RRF and\\candidate mask};
    \node [block, right=1.3cm of fuse] (rank) {Cross-encoder\\reranking};
    \node [llmblock, below=1.15cm of rank] (generate) {Grounded prompt\\gpt-4o-mini};
    \node [block, left=1.3cm of generate] (output) {Answer +\\source citations};
    \path [line] (query) -- (retrieve);
    \path [line] (query) |- (graph);
    \path [line] (retrieve) -- (fuse);
    \path [line] (graph) -- (fuse);
    \path [line] (fuse) -- (rank);
    \path [line] (rank) -- (generate);
    \path [line] (generate) -- (output);
\end{tikzpicture}%
}%
\caption{End-to-end operational workflow of L\_RAG. The graph and reranker
stages are optional branches whose effects are tested under separate frozen
ablation controls.}
\label{fig:architecture}
\end{figure*}

The component boundaries are deliberate. The retriever decides which chunks
are available, the graph can alter that set under a fixed context budget, the
reranker only reorders a candidate pool, and the generator sees a serialized
top-10 context. This separation prevents a generation metric from being
mistaken for evidence-retrieval quality.

\subsection{Offline Data and Index Construction}
\label{sec:offline-construction}

Offline construction branches from the normalized records into three linked
artifacts. The structured-text branch emits provisions and chunks with stable
parent pointers. The graph branch emits document, provision, chunk, and
relationship nodes plus validity and authority metadata. The retrieval branch
embeds either chunk-only or header-plus-chunk text into a FAISS index and
builds tokenized BM25 shards for configurations that enable sparse retrieval.

The dense index uses L2-normalized vectors and inner-product search, which is
equivalent to cosine ranking for normalized embeddings. For the completed
hybrid runs, BM25 is built as 10 independent shards with shard-local IDF and
average-document-length statistics; shard candidates are merged by BM25 score
before RRF. The graph and index artifacts keep the same chunk IDs, so a
retrieved vector row can be joined back to its document and provision without
fuzzy matching. Frozen manifests record the corpus snapshot, model
checkpoint, index path, configuration name, and benchmark version used to
create each evaluation run.

\subsection{Online Query-Processing Workflow}
\label{sec:online-query}

At query time, the system first preserves the raw Vietnamese question and any
explicit date or scope constraint. The query encoder produces an E5 vector;
the dense branch retrieves a broad seed set, while the optional BM25 branch
searches tokenized shards. The branches are merged with RRF when hybrid mode
is enabled. In the completed retrieval/graph family, a graph variant keeps
the top five seeds, expands them within their parent provisions with
\texttt{max\_hop=1}, and replaces the baseline top-10 context with at most
ten expanded chunks. This is a complete context-selection policy: it uses
broad validity filtering and does not apply a query-time validity overlay or
reranking. The resulting candidates are trimmed to the family-specific
budget, reranked only in the reranker family, and serialized as the final
context.

Insufficient evidence follows an observable fallback path. If no candidate
survives a valid scope filter, if the context contains no usable text, or if
the prompt's evidence threshold is not met, the retrieval/graph prompt asks
the model to emit the Vietnamese fallback phrase
\textit{``Không có đủ thông tin trong ngữ cảnh được cung cấp.''}; the
fixed-context reasoning runs use the visible marker
\texttt{INSUFFICIENT\_CONTEXT}. These strings are saved for evaluation and
are not treated as proof that a question was legally unanswerable. Invalid
filters and unresolved external references remain visible in the run record
instead of being silently repaired.

\subsection{Knowledge Graph Construction}
\label{sec:graph-construction}

The G-LRAG graph uses Document ($D$), Provision ($P$), Chunk ($C$), and
ExternalStub nodes. Structural edges encode
\texttt{DOCUMENT\_HAS\_PROVISION}, \texttt{PROVISION\_HAS\_CHUNK},
\texttt{CHUNK\_NEXT}, and \texttt{PROVISION\_NEXT}. Legal-relation edges encode
citations, amendments, replacements, and repeals. Each edge retains source
provenance and is available to query-time traversal only after direction and
endpoint checks.

The overlay can fold validity fields relative to an \texttt{as\_of\_date}
parameter and can sort candidates by authority rank. A simplified precedence
ordering available to the system is
\begin{equation}
\text{Constitution} \succ \text{Law/Code} \succ
\text{Decree} \succ \text{Circular}.
\end{equation}
Reading-order traversal follows adjacent chunks or provisions to recover
nearby context; hop 1 covers the seed provision and additional hops reach
subsequent provisions. In the completed graph ablation, traversal is limited
to five seeds and at most ten retained chunks, with no query-time validity
overlay or reranking. This policy can replace strong seed evidence and is
therefore evaluated as a change to the candidate set, not as a free context
addition.

\subsection{Retrieval and Ranking}
\label{sec:retrieval-ranking}

The dense encoder is \texttt{intfloat/multilingual-e5-large}. For the
metadata condition, the indexed string is
\begin{equation}
\begin{aligned}
\text{Text}_{\mathrm{embed}}={}&\text{DocumentTitle}
\;\parallel\;\text{ProvisionHeader}\\
&\;\parallel\;\text{ChunkContent}.
\end{aligned}
\end{equation}
the chunk-only condition removes the first two fields. The FAISS index uses
\texttt{IndexFlatIP}. Sparse retrieval uses \texttt{rank\_bm25} with
\texttt{underthesea} tokenization and $k_1=1.5$, $b=0.75$. Dense and sparse
lists are fused by
\begin{equation}
\mathrm{RRF}(d)=\sum_{m\in\{\mathrm{dense},\mathrm{BM25}\}}
\frac{1}{60+r_m(d)}.
\end{equation}

The dense/hybrid/graph family retrieves 50 seeds and retains at most ten final
chunks. The reranker family reconstructs a common RRF pool of 30 candidates
($C_{30}$), then retains ten ($T_{10}$). mMiniLM, Qwen3, BGE, and Jina score
the full query--candidate pair with a cross-encoder. Models are loaded,
executed, and unloaded sequentially to avoid GPU memory interference. The
reranker can improve ordering inside $C_{30}$, but it cannot recover evidence
that never entered the pool.

\subsection{Answer Generation}
\label{sec:answer-generation}

The context serializer writes each retained chunk with its chunk ID, document
title, provision header, and passage text. All generators use temperature
$0.0$, but output limits are family-specific: the retrieval/graph runs set no
explicit client-side token cap, while the fixed-context reasoning comparison
uses a 1,024-token limit. The retrieval/graph prompt asks for Vietnamese
answers and uses the fallback phrase
\textit{``Không có đủ thông tin trong ngữ cảnh được cung cấp.''} when the
evidence is insufficient; the reasoning runs use the visible marker
\texttt{INSUFFICIENT\_CONTEXT}. A citation should identify the governing
article or instrument.

The observable output record contains the final answer, citation markers,
parser fields, refusal marker, latency, and the ordered context IDs. It does
not contain hidden chain-of-thought, reasoning-token accounting, or a human
claim-entailment annotation. Consequently, the evaluation can report output
format and lexical overlap but cannot infer legal correctness from a fluent
answer or a well-formed citation.

\subsection{End-to-End Worked Example}
\label{sec:worked-example}

Consider a benchmark question asking whether a specified land-use action is
permitted under a named Vietnamese instrument. The source record is first
joined to its document ID and effective date. Its controlling article is
represented as a provision, and the provision is split into ordered chunks
whose IDs remain linked to the parent. The query encoder retrieves candidate
chunks; the graph can add the adjacent clause or a cited instrument; RRF and
the reranker order the retained candidates; and the serializer presents the
top ten with their provenance.

The generator then produces an answer in the requested format and attaches
the document and provision identifiers used as citations. If the retrieved
context contains only a related definition but not the operative condition,
the intended output is \texttt{INSUFFICIENT\_CONTEXT}. The example illustrates
why the report distinguishes document success from exact chunk recall and
citation presence from citation entailment: each is a different observable
stage in the same traceable workflow.
